% ===================================================================
%  ਸਾਰਣੀ ਨੰ. 10 — ਸਮੁੰਦਰ। — ਬੰਦਰਗਾਹ।
%  Traduction 2026 d'apres texto/io, controlee sur texto/fr, texto/en
%  et texto/hi. Consigne : voir texto/pa/10-sarni-01.tex.
%
%  Le francais decale sa numerotation d'un cran a partir de l'alinea 8.
%  C'est une coquille de l'imprimeur francais, conservee dans la
%  transcription ; le gourmoukhi suit l'ido, qui compte juste.
% ===================================================================

%%K t10-apar-1 apar
\VUsaut{4.0mm}
\VUtitre{15.0pt}{60}{ਸਾਰਣੀ ਨੰ. 10}
\VUsaut{3.0mm}
\VUpk{11.4pt}{40}{ਸਮੁੰਦਰ। --- ਬੰਦਰਗਾਹ।}
\VUsaut{3.0mm}

%%K t10-01-1 p
1. --- ਗਰਮੀਆਂ ਵਿਚ ਜਦੋਂ ਕੁਝ ਲੋਕ ਪਹਾੜਾਂ ਵਿਚ ਸ਼ਾਂਤੀ ਤੇ ਠੰਢੀ ਹਵਾ ਲੱਭਦੇ ਹਨ, ਦੂਜੇ
ਸਮੁੰਦਰ ਦੇ ਕੰਢੇ ਜਾਣਾ ਵੱਧ ਪਸੰਦ ਕਰਦੇ ਹਨ। ਪਿਛਲੇ ਵਰ੍ਹੇ ਅਸੀਂ ਇਹੋ ਕੀਤਾ, ਕਿਉਂਕਿ ਸਾਨੂੰ
\VUgras{ਮਹਾਂਸਾਗਰ} \textsuperscript{(1)} ਬਹੁਤ ਚੰਗਾ ਲਗਦਾ ਹੈ।

%%K t10-02-1 p
2. --- ਮੈਨੂੰ ਤਾਂ \VUgras{ਦਿਸਹੱਦੇ} \textsuperscript{(2)} ਤਕ ਫੈਲੀ ਉਸ ਵੱਡੀ ਪਾਣੀ ਦੀ
ਚਾਦਰ ਨੂੰ ਤੱਕਣ ਵਿਚ ਬੜਾ ਸੁਖ ਮਿਲਦਾ ਹੈ, ਤੇ ਇਹ ਵੇਖਣ ਵਿਚ ਕਿ ਵੱਡੇ \VUgras{ਭਾਫ਼ ਦੇ ਜਹਾਜ਼}
\textsuperscript{(3)} ਅਮਰੀਕਾ ਜਾਂਦੇ ਮੁਸਾਫ਼ਰਾਂ ਨੂੰ ਲੈ ਕੇ ਲੰਮੇ \VUgras{ਸਮੁੰਦਰੀ ਸਫ਼ਰ}
ਉੱਤੇ ਨਿਕਲਦੇ ਹਨ।

%%K t10-03-1 p
3. --- ਸੋ ਮੇਰੇ ਪਿਤਾ, ਜੋ ਜਾਣਦੇ ਹਨ ਕਿ ਮੈਨੂੰ ਕੀ ਚੰਗਾ ਲਗਦਾ ਹੈ, ਛੁੱਟੀਆਂ ਲਈ ਹਮੇਸ਼ਾ
ਸਾਡੀ ਕਿਸੇ ਵੱਡੀ \VUgras{ਜਲ-ਸੈਨਾ ਦੀ ਬੰਦਰਗਾਹ} ਕੋਲ ਦਾ ਕੋਈ ਬਹੁਤ ਸ਼ਾਂਤ ਕੰਢਾ ਚੁਣਦੇ ਹਨ।

%%K t10-03-2 p
ਸਾਡੇ ਪਿਛਲੇ ਟਿਕਾਅ ਵਿਚ \VUgras{ਬੇੜਾ} ਆ ਕੇ ਉਸ ਵੱਡੀ \VUgras{ਲਹਿਰ-ਰੋਕ}
\textsuperscript{(4)} ਦੇ ਪਿੱਛੇ ਲੰਗਰ ਸੁੱਟ ਗਿਆ ਜੋ \VUgras{ਜਲ-ਸੈਨਾ ਦੀ ਬੰਦਰਗਾਹ}
\textsuperscript{(5)} ਨੂੰ ਬੰਦ ਕਰਦੀ ਤੇ ਉਸ ਦੀ ਰਾਖੀ ਕਰਦੀ ਹੈ, ਤੇ ਮੈਂ ਸਕੂਨ ਨਾਲ ਵੱਖੋ-ਵੱਖ
\VUgras{ਜੰਗੀ ਜਹਾਜ਼} ਤੱਕ ਸਕਿਆ — ਨਿੱਕੀ \VUgras{ਪਣਡੁੱਬੀ} \textsuperscript{(10)}
ਤੋਂ, ਜੋ ਚੁੱਭੀ ਮਾਰ ਕੇ ਪਾਣੀ ਦੇ ਹੇਠਾਂ ਗੁੰਮ ਹੋ ਜਾਂਦੀ ਹੈ, ਵੱਡੇ \VUgras{ਜੰਗੀ ਜਹਾਜ਼}
\textsuperscript{(9)} ਤਕ, ਜਿਸ ਨੂੰ ਹੁਣ ਤਕ \VUgras{ਤਾਰਪੀਡੋ ਕਿਸ਼ਤੀ}
\textsuperscript{(8)} ਦੇ ਹੱਲਿਆਂ ਤੋਂ ਬਿਨਾਂ ਕੁਝ ਡਰਾਉਂਦਾ ਨਾ ਸੀ।

%%K t10-tit-2 sub
\VUsaut{3.0mm}
\VUpk{10.2pt}{40}{ਨਿੱਕੀਆਂ ਕਿਸ਼ਤੀਆਂ।}
\VUsaut{3.0mm}

%%K t10-04-1 p
4. --- ਉਹ ਪਹਿਲਾਂ ਹੀ \VUgras{ਸਿਆਣਪ} ਨਾਲ ਮੋਟੇ ਧਾਤ ਦੇ ਕਵਚ ਦੀ ਕਮਜ਼ੋਰ ਥਾਂ ਲੱਭ ਲੈਂਦੀ
ਸੀ ਤੇ ਉੱਥੇ ਉਹ ਖ਼ਤਰਨਾਕ \VUgras{ਤਾਰਪੀਡੋ} ਗੱਡ ਦਿੰਦੀ ਸੀ ਜਿਸ ਦੇ ਫਟਣ ਨਾਲ ਜਹਾਜ਼ ਡੁੱਬ
ਜਾਂਦਾ ਹੈ; ਪਰ ਉਹ ਪਣਡੁੱਬੀ ਤੋਂ ਫਿਰ ਵੀ ਘੱਟ ਡਰਾਉਣੀ ਸੀ, ਕਿਉਂਕਿ ਤਬਾਹਕੁੰਨ ਉਸ ਦਾ ਪਿੱਛਾ
ਸੌਖਾ ਹੀ ਕਰ ਸਕਦੇ ਹਨ।

%%K t10-05-1 p
5. --- ਮੈਂ ਤੇਜ਼ ਕਵਚਬੰਦ \VUgras{ਕਰੂਜ਼ਰ} \textsuperscript{(6)} ਵੀ ਵੇਖ ਸਕਿਆ, ਜੋ ਸੱਚੇ
ਜੰਗੀ ਜਹਾਜ਼ਾਂ ਜਿੰਨੇ ਹੀ ਅਭੇਦ ਹਨ, ਕਈ \VUgras{ਫ਼ੌਜੀ ਜਹਾਜ਼} \textsuperscript{(12)},
ਤਿੰਨ-ਚਾਰ \VUgras{ਤੋਪ-ਕਿਸ਼ਤੀਆਂ} \textsuperscript{(7)}, ਤੇ ਅੰਤ ਵਿਚ ਇਕ
\VUgras{ਫ਼ਰੀਗੇਟ} \textsuperscript{(11)}।
\VUgras{ਉਸ ਬੇੜੇ ਨੂੰ ਲੰਗਰ ਚੁੱਕਦੇ ਵੇਖਣਾ ਸਭ ਤੋਂ ਦਿਲਚਸਪ ਹੈ।}

%%K t10-06-1 p
6. --- ਫ਼ੌਲਾਦ ਦੇ ਉਨ੍ਹਾਂ ਵੱਡੇ ਢੇਰਾਂ ਤੇ ਉਨ੍ਹਾਂ ਸਜੀਲੇ \VUgras{ਤਿੰਨ-ਮਸਤੂਲੀ} ਜਾਂ
ਸੋਹਣੇ \VUgras{ਬਾਦਬਾਨੀ ਜਹਾਜ਼ਾਂ} \textsuperscript{(13)} ਦੀ ਬਣਤਰ ਵਿਚ ਕਿੰਨਾ ਫ਼ਰਕ ਹੈ, ਜੋ
ਹੁਣ ਵੀ ਵਪਾਰੀ ਬੇੜੇ ਵਿਚ ਕੰਮ ਆਉਂਦੇ ਹਨ ਤੇ ਜਿਨ੍ਹਾਂ ਨੂੰ ਮੈਂ \VUgras{ਘਾਟ-ਬੰਨ੍ਹ}
\textsuperscript{(14)} ਉੱਤੇ ਟਹਿਲਦੇ ਹੋਏ ਪੂਰੇ ਬਾਦਬਾਨ ਖੋਲ੍ਹ ਕੇ ਬੰਦਰਗਾਹ ਵਿਚ ਆਉਂਦੇ
ਵੇਖਦਾ ਸੀ। ਇਹ ਸੱਚ ਹੈ ਕਿ \VUgras{ਹਵਾ} ਉਲਟ ਹੋਣ ਉੱਤੇ ਉਹ ਆਪਣਾ ਬਹੁਤ ਫ਼ਾਇਦਾ ਗੁਆ ਦਿੰਦੇ
ਸਨ। ਗੋਦੀ ਵਿਚ ਵੜਨ ਲਈ ਉਨ੍ਹਾਂ ਨੂੰ ਸਾਰੇ ਬਾਦਬਾਨ ਸਮੇਟਣੇ ਪੈਂਦੇ ਸਨ, ਤੇ ਇਕ
\VUgras{ਖਿੱਚਣ ਵਾਲੀ ਕਿਸ਼ਤੀ} \textsuperscript{(16)} ਚਾਹੀਦੀ ਹੁੰਦੀ ਸੀ ਜੋ ਉਨ੍ਹਾਂ ਨੂੰ ਲੈਣ
ਜਾਵੇ ਤੇ ਸਧਾਰਨ \VUgras{ਮਾਲ ਢੋਣ ਵਾਲੀਆਂ ਕਿਸ਼ਤੀਆਂ} \textsuperscript{(17)} ਵਾਂਗ ਘਾਟ ਤਕ ਲੈ
ਆਵੇ।

%%K t10-tit-3 sub
\VUsaut{3.0mm}
\VUpk{10.2pt}{40}{ਵਪਾਰੀ ਬੰਦਰਗਾਹ।}
\VUsaut{3.0mm}

%%K t10-07-1 p
7. --- ਜਿਸ ਬੰਦਰਗਾਹ ਦੀ ਮੈਂ ਗੱਲ ਕਰ ਰਿਹਾ ਹਾਂ ਉਹ ਇਸੇ ਲਈ ਇੰਨੀ ਦਿਲਚਸਪ ਹੈ ਕਿ ਉਸ ਵਿਚ
ਸਿਰਫ਼ ਵੱਡੇ ਕੰਮਾਂ ਨਾਲ ਬਣਾਈ ਹੋਈ ਜਲ-ਸੈਨਾ ਦੀ ਬੰਦਰਗਾਹ ਹੀ ਨਹੀਂ, ਸਗੋਂ ਇਕ ਵੱਡੇ ਦਰਿਆ ਦੇ
\VUgras{ਮੁਹਾਣੇ} ਉੱਤੇ ਇਕ ਹੈਰਾਨ ਕਰ ਦੇਣ ਵਾਲੀ \VUgras{ਵਪਾਰੀ ਬੰਦਰਗਾਹ} ਵੀ ਹੈ।

%%K t10-08-1 p
8. --- ਦੋਹਾਂ ਗੋਦੀਆਂ ਨੂੰ ਵੱਖ ਕਰਦੀ ਜ਼ਮੀਨ ਦੀ ਪੱਟੀ ਕਿਲ੍ਹਾਬੰਦ ਹੈ ਤੇ ਸਮੁੰਦਰ ਤਕ ਫੈਲੀ
\VUgras{ਤੋਪਾਂ ਦੀ ਕਤਾਰ} ਤੇ \VUgras{ਬੁਰਜਾਂ} ਨਾਲ ਰੱਖੀ ਹੋਈ ਹੈ। \VUgras{ਫ਼ਸੀਲਾਂ}
\textsuperscript{(15)} ਉੱਤੇ ਟਹਿਲਣ ਲਈ ਖ਼ਾਸ ਆਗਿਆ ਚਾਹੀਦੀ ਹੈ। ਮੈਨੂੰ ਉਹ ਆਪਣੇ ਚਾਚਾ ਦੇ
ਜ਼ਰੀਏ ਸੌਖੀ ਹੀ ਮਿਲ ਗਈ ਸੀ, ਜੋ \VUgras{ਜਹਾਜ਼-ਸਾਜ਼ੀ ਦੇ ਕਾਰਖ਼ਾਨੇ} \textsuperscript{(18)}
ਵਿਚ ਇੰਜੀਨੀਅਰ ਹਨ, ਤੇ ਮੈਂ ਵੱਡੇ ਛੱਪਰਾਂ ਦੇ ਹੇਠਾਂ ਬਣਦੇ ਜਹਾਜ਼ ਵੇਖਣ ਗਿਆ।

%%K t10-09-1 p
9. --- ਮੈਂ ਇਕ ਜੰਗੀ ਜਹਾਜ਼ ਨੂੰ ਪਾਣੀ ਵਿਚ ਉਤਰਦੇ ਵੀ ਵੇਖਿਆ, ਤੇ ਮੈਨੂੰ ਉਹ ਰੋਮਾਂਚ ਯਾਦ ਹੈ
ਜੋ ਸਾਰੀ ਭੀੜ ਨੇ ਮਹਿਸੂਸ ਕੀਤਾ ਜਦੋਂ ਉਹ ਵੱਡਾ ਢੇਰ, ਆਪਣੇ ਭਰੋਸੇ ਛੱਡਿਆ ਗਿਆ, ਆਪਣੀ ਪਾਲਣੀ
ਉੱਤੇ ਸਰਕਣ ਲੱਗਾ ਤੇ ਦਰਿਆ ਦੇ ਪਾਣੀ ਉੱਤੇ ਤਰਨ ਆ ਗਿਆ।

%%K t10-10-1 p
10. --- ਕਾਰਖ਼ਾਨੇ ਤਕ ਪਹੁੰਚਣ ਲਈ \VUgras{ਪਰੇਡ-ਮੈਦਾਨ} \textsuperscript{(19)} ਲੰਘਣਾ
ਪੈਂਦਾ ਹੈ, ਜਿੱਥੇ ਛਾਉਣੀ ਦੀ ਫ਼ੌਜ ਹਰ ਦਿਨ ਕਵਾਇਦ ਨੂੰ ਆਉਂਦੀ ਹੈ, ਤੇ ਲੋਹੇ ਦੀ ਉਹ
\VUgras{ਪੁਲੀ} \textsuperscript{(20)} ਲੰਘਣੀ ਪੈਂਦੀ ਹੈ ਜੋ ਪਰੇਡ-ਮੈਦਾਨ ਨੂੰ
\VUgras{ਅਸਲਾਖ਼ਾਨੇ} \textsuperscript{(21)} ਨਾਲ ਜੋੜਦੀ ਹੈ।

%%K t10-tit-4 sub
\VUsaut{3.0mm}
\VUpk{10.2pt}{40}{ਅਸਲਾਖ਼ਾਨਾ ਤੇ ਗੋਦੀਆਂ।}
\VUsaut{3.0mm}

%%K t10-11-1 p
11. --- ਆਪਣੇ ਚਾਚਾ ਨਾਲ ਮੈਂ ਉਹ ਵੱਡਾ ਅਦਾਰਾ ਵੇਖਿਆ, ਜੋ \VUgras{ਤੋਪਾਂ}
\textsuperscript{(22)}, \VUgras{ਗੋਲਿਆਂ} \textsuperscript{(23)} ਤੇ \VUgras{ਖੋਲਾਂ} ਨਾਲ
ਭਰਿਆ ਹੈ। ਮੈਂ ਉਹ \VUgras{ਲੋਹ-ਸ਼ਾਲਾ} \textsuperscript{(24)} ਵੀ ਵੇਖੀ ਜਿੱਥੇ ਭਾਫ਼ ਦੇ
ਜਹਾਜ਼ਾਂ ਦੇ \VUgras{ਬੌਇਲਰ} \textsuperscript{(25)} ਬਣਦੇ ਹਨ।

%%K t10-12-1 p
12. --- ਬਹੁਤ ਕੋਲ ਹੀ \VUgras{ਗੋਦੀਆਂ} \textsuperscript{(26)} ਹਨ — ਸੁੱਕੀਆਂ ਗੋਦੀਆਂ —
ਤੇ ਉਹ ਬਹੁਤ ਜ਼ਿਕਰ ਜੋਗ \VUgras{ਤਰਦੀਆਂ ਗੋਦੀਆਂ} \textsuperscript{(27)}, ਜਿਨ੍ਹਾਂ ਵਿਚ ਮੈਂ
ਮੁਰੰਮਤ ਹੁੰਦੇ ਇਕ ਜਹਾਜ਼ ਉੱਤੇ ਕੋਲੋਂ ਕਿਸੇ ਜਹਾਜ਼ ਦਾ \VUgras{ਪੱਖਾ} \textsuperscript{(28)},
\VUgras{ਪੇਂਦਾ} \textsuperscript{(29)} ਤੇ \VUgras{ਪਤਵਾਰ} \textsuperscript{(30)} ਵੇਖ
ਸਕਿਆ।

%%K t10-13-1 p
13. --- ਵਪਾਰੀ ਬੰਦਰਗਾਹ ਜਲ-ਸੈਨਾ ਦੀ ਬੰਦਰਗਾਹ ਜਿੰਨੀ ਹੀ ਵੇਖਣ ਜੋਗ ਹੈ, ਤੇ ਮੈਂ ਕਈ ਘੰਟੇ
ਉਨ੍ਹਾਂ ਘਾਟਾਂ ਉੱਤੇ ਲੰਘਾਏ ਜਿੱਥੇ ਦਰਿਆ ਚੜ੍ਹਨ ਵਾਲੇ \VUgras{ਚੱਕਰ ਵਾਲੇ ਭਾਫ਼ ਜਹਾਜ਼}
\textsuperscript{(31)} ਬੰਨ੍ਹੇ ਰਹਿੰਦੇ ਹਨ, ਤੇ ਉਸ \VUgras{ਤਿਪਾਈ ਕਰੇਨ}
\textsuperscript{(32)} ਨੂੰ ਕੰਮ ਕਰਦੇ ਵੇਖਦੇ ਹੋਏ, ਜੋ ਵੱਡੇ ਬੋਝ ਖੰਭਾਂ ਵਾਂਗ ਚੁੱਕ ਲੈਂਦੀ ਹੈ।

%%K t10-tit-5 sub
\VUsaut{4.0mm}
\VUpk{10.2pt}{40}{ਰਾਹ-ਦਸੇਰਾ ਮਲਾਹ।}
\VUsaut{3.0mm}

%%K t10-14-1 p
14. --- ਮੈਂ ਉਨ੍ਹਾਂ \VUgras{ਅੰਤਰ-ਮਹਾਂਸਾਗਰੀ ਜਹਾਜ਼ਾਂ} \textsuperscript{(34)} ਨੂੰ ਤੱਕਿਆ
ਜੋ ਆਪਣੀਆਂ \VUgras{ਝੰਡੀਆਂ} \textsuperscript{(35)} ਲਹਿਰਾਉਂਦੇ ਬੰਦਰਗਾਹ ਛੱਡਦੇ ਹਨ,
ਜਿਨ੍ਹਾਂ ਨੂੰ ਇਕ ਸਿਆਣਾ \VUgras{ਰਾਹ-ਦਸੇਰਾ ਮਲਾਹ} \textsuperscript{(36)} ਚਲਾਉਂਦਾ ਹੈ, ਜੋ
\VUgras{ਬੋਇਆਂ} \textsuperscript{(40)} ਤੇ \VUgras{ਇਸ਼ਾਰੇ ਦੇ ਥੰਮ੍ਹਾਂ} ਨਾਲ ਨਿਸ਼ਾਨ ਲਾਏ
\VUgras{ਪਾਣੀ ਦੇ ਰਾਹ} ਉੱਤੇ ਹੈਰਾਨ ਕਰ ਦੇਣ ਵਾਲੇ ਢੰਗ ਨਾਲ ਤੁਰਨਾ ਜਾਣਦਾ ਹੈ, ਤੇ ਉਨ੍ਹਾਂ
\VUgras{ਮਾਲ ਕਿਸ਼ਤੀਆਂ} \textsuperscript{(41)} ਤੋਂ ਬਚਦਾ ਹੈ ਜੋ ਆਪਣੇ ਇਕੱਲੇ ਚੌਕੋਰ ਬਾਦਬਾਨ
ਨਾਲ ਬੜੀ ਔਖ ਨਾਲ ਚਲਾਈਆਂ ਜਾਂਦੀਆਂ ਹਨ। ਮੈਂ \VUgras{ਅਗਲੇ ਹਿੱਸੇ} \textsuperscript{(37)}
ਉੱਤੇ ਦੂਜੇ ਦਰਜੇ ਦੀਆਂ \VUgras{ਕੋਠੜੀਆਂ} \textsuperscript{(39)} ਪਛਾਣ ਸਕਿਆ, ਤੇ
\VUgras{ਪਿਛਲੇ ਹਿੱਸੇ} \textsuperscript{(38)} ਉੱਤੇ ਪਹਿਲੇ ਦਰਜੇ ਦੇ ਮੁਸਾਫ਼ਰਾਂ ਦੇ ਕਮਰੇ।

%%K t10-15-1 p
15. --- ਕਦੇ ਕਦੇ ਮੈਂ ਕਿਸੇ \VUgras{ਮਾਲ-ਬਰਦਾਰ ਜਹਾਜ਼} \textsuperscript{(42)} ਦੀ ਲਦਾਈ ਵੀ
ਵੇਖਦਾ, ਜਿਸ ਵਿਚ \VUgras{ਗੁਦਾਮ} \textsuperscript{(43)} ਤੇ
\VUgras{ਬੰਦਰਗਾਹ ਦੇ ਸਟੇਸ਼ਨ} \textsuperscript{(44)} ਦਾ ਮਾਲ ਠੋਸਿਆ ਜਾਂਦਾ, ਜੋ
\VUgras{ਘਾਟ ਦੀ ਰੇਲ-ਲਾਈਨ} \textsuperscript{(45)} ਦੀਆਂ ਕਈ ਗੱਡੀਆਂ ਲਿਆਉਂਦੀਆਂ ਸਨ।

%%K t10-tit-6 sub
\VUsaut{3.0mm}
\VUpk{10.2pt}{40}{ਕਿਸ਼ਤੀ ਦੀ ਸੈਰ।}
\VUsaut{3.0mm}

%%K t10-16-1 p
16. --- ਮੈਂ ਅਕਸਰ \VUgras{ਤਰਦੀ ਗੋਦੀ} \textsuperscript{(53)} ਵਿਚ ਕਿਸ਼ਤੀ ਚਲਾਉਂਦਾ ਸੀ,
ਜਿੱਥੇ ਨਿੱਕੀਆਂ \VUgras{ਕਿਸ਼ਤੀਆਂ} \textsuperscript{(46)} ਪਨਾਹ ਲੈਂਦੀਆਂ ਹਨ, ਤੇ
\VUgras{ਚੱਪੂ} \textsuperscript{(47)} ਚਲਾਉਣਾ ਮੇਰੇ ਲਈ ਆਨੰਦ ਸੀ। ਮੈਂ \VUgras{ਪਟੜੇ}
\textsuperscript{(48)} ਉੱਤੇ ਚੜ੍ਹ ਜਾਂਦਾ, ਕਿਸੇ \VUgras{ਬਾਦਬਾਨੀ ਕਿਸ਼ਤੀ}
\textsuperscript{(49)} ਦੇ, \VUgras{ਰੱਸੀਆਂ} \textsuperscript{(52)} ਦੇ ਸਹਾਰੇ
\VUgras{ਮਸਤੂਲ} \textsuperscript{(51)} ਉੱਤੇ \VUgras{ਆੜੇ ਡੰਡਿਆਂ}
\textsuperscript{(50)} ਤਕ ਚੜ੍ਹਦਾ, ਫਿਰ ਇਕ \VUgras{ਡੂੰਗੀ} \textsuperscript{(54)} ਵਿਚ
ਆਪਣੀ ਸੈਰ ਜਾਰੀ ਰੱਖਦਾ, ਭਾਰੀਆਂ \VUgras{ਮੱਛੀ ਦੀਆਂ ਕਿਸ਼ਤੀਆਂ} \textsuperscript{(55)}
ਵਿਚਕਾਰ, ਜਿੱਥੇ \VUgras{ਜਾਲ} \textsuperscript{(56)} ਸੁੱਕਣ ਨੂੰ ਟੰਗੇ ਰਹਿੰਦੇ ਹਨ।

%%K t10-17-1 p
17. --- \VUgras{ਘਾਟ} \textsuperscript{(58)} ਉੱਤੇ ਮੇਰੇ ਸਾਹਮਣੇ ਤੋਂ ਕਈ ਤਰ੍ਹਾਂ ਦਾ
\VUgras{ਮਾਲ} \textsuperscript{(59)} ਲੰਘਦਾ, \VUgras{ਸੰਦੂਕਾਂ} \textsuperscript{(60)}
ਜਾਂ \VUgras{ਗੱਠੜੀਆਂ} \textsuperscript{(61)} ਵਿਚ ਬੰਨ੍ਹਿਆ। ਉਹ ਮਾਲ ਕਿਸ਼ਤੀਆਂ ਦੀ
\VUgras{ਲਦਾਈ} \textsuperscript{(63)} ਬਣਾਉਂਦੇ ਸਨ, ਤੇ ਤਕੜੀਆਂ \VUgras{ਕਰੇਨਾਂ}
\textsuperscript{(62)} ਉਨ੍ਹਾਂ ਨੂੰ ਚੁੱਕ ਕੇ \VUgras{ਟਰੱਕਾਂ} \textsuperscript{(64)} ਉੱਤੇ
ਰੱਖਦੀਆਂ ਸਨ, ਜਦੋਂ ਤਕ \VUgras{ਢੋਣ ਵਾਲੇ} \VUgras{ਚੁੰਗੀ-ਘਰ} \textsuperscript{(65)} ਵਿਚ
ਉਨ੍ਹਾਂ ਦਾ ਐਲਾਨ ਕਰਦੇ ਤੇ ਮਹਿਸੂਲ ਤਾਰਦੇ।

%%K t10-18-1 p
18. --- ਅੰਤ ਵਿਚ, \VUgras{ਘੁੰਮਣ ਵਾਲੇ ਪੁਲ} \textsuperscript{(66)} ਜਾਂ
\VUgras{ਜਲ-ਦਰਵਾਜ਼ੇ} \textsuperscript{(69)} ਤਕ ਪਹੁੰਚ ਕੇ, ਉਸ \VUgras{ਪੁੱਟਣ ਵਾਲੀ
ਕਿਸ਼ਤੀ} \textsuperscript{(68)} ਕੋਲੋਂ ਲੰਘਦੇ ਹੋਏ ਜੋ \VUgras{ਨਹਿਰ}
\textsuperscript{(67)} ਸਾਫ਼ ਕਰਦੀ ਹੈ, ਮੈਂ \VUgras{ਇਸ਼ਾਰੇ ਦੇ ਮਸਤੂਲ}
\textsuperscript{(70)} ਕੋਲ ਘਾਟ ਉੱਤੇ ਉਤਰ ਜਾਂਦਾ।

%%K t10-19-1 p
19. --- ਉਸੇ ਘਾਟ ਦੇ ਦੂਜੇ ਪਾਸੇ ਉਹ \VUgras{ਨਿੱਕੀਆਂ ਕਿਸ਼ਤੀਆਂ} \textsuperscript{(71)}
ਲਗਦੀਆਂ ਹਨ ਜੋ ਬੇੜੇ ਦੇ ਅਫ਼ਸਰਾਂ ਨੂੰ ਕੰਢੇ ਲਿਆਉਂਦੀਆਂ ਹਨ। ਮਲਾਹਾਂ ਨੂੰ ਤਾਲ ਨਾਲ ਚੱਪੂ
ਚੁੱਕਦੇ ਵੇਖਣਾ ਸੋਹਣਾ ਦ੍ਰਿਸ਼ ਹੈ, ਜਦੋਂ \VUgras{ਛੱਲ} \textsuperscript{(72)} ਉੱਤੇ ਚੜ੍ਹੀ
ਕਿਸ਼ਤੀ ਸੌਖੀ ਹੀ ਉਤਰਨ ਦੀਆਂ ਪੌੜੀਆਂ ਤਕ ਆ ਲਗਦੀ ਹੈ। ਇੱਥੋਂ ਹੀ \VUgras{ਜਵਾਰ} ਤੇ
\VUgras{ਕੰਢੇ} \textsuperscript{(73)} ਉੱਤੇ \VUgras{ਭਾਟੇ} ਤੇ \VUgras{ਜਵਾਰ} ਦੀ ਚਾਲ ਸਭ
ਤੋਂ ਚੰਗੀ ਵੇਖੀ ਜਾ ਸਕਦੀ ਹੈ।

%%K t10-tit-7 sub
\VUsaut{3.0mm}
\VUpk{10.2pt}{40}{ਅਫ਼ਸਰਾਂ ਦਾ ਕਲੱਬ।}
\VUsaut{3.0mm}

%%K t10-20-1 p
20. --- ਘਰ ਮੁੜਨ ਲਈ ਮੈਂ \VUgras{ਬਿਜਲੀ ਦੀ ਟਰਾਮ} \textsuperscript{(74)} ਲੈਂਦਾ, ਜਿਸ ਦਾ
\VUgras{ਟਿਕਾਣਾ} \textsuperscript{(75)} ਅਫ਼ਸਰਾਂ ਦੇ ਕਲੱਬ ਦੇ ਸਾਹਮਣੇ ਖੜੇ \VUgras{ਥੰਮ੍ਹ}
ਉੱਤੇ ਹੈ।

%%K t10-20-2 p
ਕਲੱਬ ਦੇ \VUgras{ਬਾਲਾਖ਼ਾਨੇ} \textsuperscript{(76)} ਤੋਂ, ਜੋ \VUgras{ਲੰਗਰਾਂ}
\textsuperscript{(77)}, ਤੋਪਾਂ ਤੇ ਗੋਲਿਆਂ ਨਾਲ ਸਜਿਆ ਹੈ, ਮੈਂ ਅਕਸਰ ਜਲ-ਸੈਨਾ ਦੇ ਅਫ਼ਸਰਾਂ ਦਾ
ਸ਼ਾਨਦਾਰ ਜਮਘਟਾ ਵੇਖਿਆ : ਆਪਣੀਆਂ ਤਲਵਾਰਾਂ ਸਮੇਤ \VUgras{ਲੈਫ਼ਟੀਨੈਂਟ}
\textsuperscript{(78)}, \VUgras{ਤੋਪਖ਼ਾਨੇ} \textsuperscript{(79)} ਤੇ
\VUgras{ਪੈਦਲ ਫ਼ੌਜ} \textsuperscript{(80)} ਦੇ \VUgras{ਕਪਤਾਨ} ਆਪਣੀਆਂ \VUgras{ਕਿਰਚਾਂ}
ਸਮੇਤ। ਉਨ੍ਹਾਂ ਦੇ ਪਿੱਛੇ ਇਕ \VUgras{ਮਲਾਹ} \textsuperscript{(81)} ਖੜਾ ਰਹਿੰਦਾ, ਜੋ
ਉਨ੍ਹਾਂ ਨੂੰ ਇਕ \VUgras{ਦੂਰਬੀਨ} ਲਿਆ ਕੇ ਦਿੰਦਾ ਜਦੋਂ ਉਹ ਦੂਰ ਦੀ ਗੱਲ ਪਛਾਣਨਾ ਚਾਹੁੰਦੇ।

%%K t10-21-1 p
21. --- ਮੈਂ ਜਵਾਨ \VUgras{ਨਾਇਬ ਲੈਫ਼ਟੀਨੈਂਟ} \textsuperscript{(82)} ਵੀ ਵੇਖੇ, ਜੋ ਆਪਣੇ
\VUgras{ਕਮਾਂਡਰ} \textsuperscript{(83)} ਜਾਂ \VUgras{ਐਡਮਿਰਲ} \textsuperscript{(84)}
ਕੋਲੋਂ ਹੁਕਮ ਲੈਂਦੇ ਸਨ, ਜਦੋਂ ਤਕ ਇਕ \VUgras{ਕੁਆਰਟਰ ਮਾਸਟਰ} \textsuperscript{(85)}
ਉਨ੍ਹਾਂ ਨੂੰ ਜਹਾਜ਼ ਤਕ ਲੈ ਜਾਣ ਦੀ ਉਡੀਕ ਕਰਦਾ।

%%K t10-22-1 p
22. --- ਉਸ ਬਾਲਾਖ਼ਾਨੇ ਤੋਂ ਦ੍ਰਿਸ਼ ਸ਼ਾਨਦਾਰ ਹੈ : ਉੱਥੋਂ ਸਾਰਾ ਕੰਢਾ ਵਿਖਾਈ ਦਿੰਦਾ ਹੈ, ਆਪਣੇ
\VUgras{ਤੰਬੂਆਂ} \textsuperscript{(86)} ਤੇ \VUgras{ਨ੍ਹਾਉਣ ਦੀਆਂ ਝੁੱਗੀਆਂ}
\textsuperscript{(87)} ਸਮੇਤ, ਤੇ ਨ੍ਹਾਉਣ ਦੇ ਵੇਲੇ \VUgras{ਨ੍ਹਾਉਣ ਵਾਲਿਆਂ}
\textsuperscript{(88)} ਨੂੰ \VUgras{ਜੂਆ-ਘਰ} \textsuperscript{(89)} ਦੇ ਸਾਹਮਣੇ ਖ਼ੁਸ਼
ਟੱਪੂਸੀਆਂ ਮਾਰਦੇ ਵੇਖਿਆ ਜਾ ਸਕਦਾ ਹੈ।

%%K t10-23-1 p
23. --- ਕੰਢੇ ਦੇ ਸਿਰੇ ਉੱਤੇ ਕਿਨਾਰਾ ਇਕ ਨਿੱਕਾ ਚਟਾਨੀ \VUgras{ਪ੍ਰਾਇਦੀਪ}
\textsuperscript{(91)} ਬਣਾਉਂਦਾ ਹੈ, ਜਿੱਥੇ \VUgras{ਛੱਲਾਂ} \textsuperscript{(92)}
ਟਕਰਾ ਕੇ ਟੁੱਟਦੀਆਂ ਹਨ ਤੇ ਜਿਸ ਉੱਤੇ ਇਕ \VUgras{ਰੌਸ਼ਨੀ ਦਾ ਮੁਨਾਰਾ}
\textsuperscript{(93)} ਬਣਿਆ ਹੈ, ਜੋ ਰਾਤ ਨੂੰ ਮਲਾਹਾਂ ਨੂੰ ਰਾਹ ਵਿਖਾਉਂਦਾ ਹੈ। ਇਹ ਪ੍ਰਾਇਦੀਪ
ਜ਼ਮੀਨ ਨਾਲ ਸਿਰਫ਼ ਇਕ ਬਹੁਤ ਭੀੜੇ \VUgras{ਡਮਰੂ-ਮੱਧ} \textsuperscript{(90)} ਨਾਲ ਜੁੜਿਆ ਹੈ।

%%K t10-tit-8 sub
\VUsaut{3.0mm}
\VUpk{10.2pt}{40}{ਕੰਢੇ ਦਾ ਆਮ ਦ੍ਰਿਸ਼।}
\VUsaut{3.0mm}

%%K t10-24-1 p
24. --- \VUgras{ਕਗਾਰ} \textsuperscript{(94)} ਅੱਗੇ ਚਲਦਾ ਜਾਂਦਾ ਹੈ, ਸਮੁੰਦਰ ਤੋਂ ਕੱਟਿਆ
ਹੋਇਆ, ਜੋ ਉਸ ਵਿਚੋਂ \VUgras{ਖਾੜੀਆਂ} \textsuperscript{(95)} ਤੇ \VUgras{ਅੰਤਰੀਪਾਂ} ਦੀ
ਇਕ ਅਨੋਖੀ ਲੜੀ ਤਰਾਸ਼ਦਾ ਹੈ, ਤੇ ਉਸ ਨੂੰ ਬਹੁਤ ਸੋਹਣਾ ਬਣਾ ਦਿੰਦਾ ਹੈ।

%%K t10-24-2 p
\VUgras{ਪਠਾਰ} \textsuperscript{(96)} ਉੱਤੇ, ਜੋ \VUgras{ਸਮੁੰਦਰੀ ਪੱਤਣ}
\textsuperscript{(98)} ਉੱਤੇ ਝੁਕਿਆ ਹੈ, ਇਕ ਬਹੁਤ ਅਹਿਮ \VUgras{ਗੜ੍ਹੀ}
\textsuperscript{(97)} ਹੈ ਤੇ ਕੁਝ ਕੋਠੀਆਂ।

%%K t10-25-1 p
25. --- ਉਹ ਸਮੁੰਦਰੀ ਪੱਤਣ ਵੱਡੀ ਜ਼ਮੀਨ ਤੋਂ ਇਕ ਬਹੁਤ ਖ਼ਤਰਨਾਕ \VUgras{ਟਾਪੂ}
\textsuperscript{(99)} ਨੂੰ ਵੱਖ ਕਰਦਾ ਹੈ, ਜੋ \VUgras{ਚਟਾਨਾਂ} \textsuperscript{(100)}
ਤੇ \VUgras{ਭਿੱਤਾਂ} ਨਾਲ ਘਿਰਿਆ ਹੈ, ਜਿੱਥੇ ਕਿਸ਼ਤੀਆਂ ਤੇ ਵੱਡੇ ਜਹਾਜ਼ ਤਕ ਕਦੇ ਕਦੇ ਟੁੱਟ
ਜਾਂਦੇ ਹਨ। ਬਚਾਅ ਦੀ ਸਾਰੀ ਫੁਰਤੀ ਤੇ \VUgras{ਬਚਾਅ ਦੀਆਂ ਕਿਸ਼ਤੀਆਂ} ਦੇ ਛੇਤੀ ਪਹੁੰਚਣ ਦੇ
ਬਾਵਜੂਦ ਇਹ \VUgras{ਜਹਾਜ਼-ਭੰਨ} ਭਿਆਨਕ ਹੁੰਦੇ ਹਨ, ਤੇ ਸਮ-ਰਾਤ ਦੀਆਂ \VUgras{ਹਨੇਰੀਆਂ}
ਵਿਚ ਕਈ ਜਹਾਜ਼ ਆਪਣੇ ਸਾਰੇ ਬੰਦਿਆਂ ਤੇ ਮਾਲ ਸਮੇਤ ਡੁੱਬ ਚੁੱਕੇ ਹਨ।

%%K t10-25-2 p
ਇਸ ਗੁਆਂਢ ਦੇ ਬਾਵਜੂਦ \VUgras{ਬੰਦਰਗਾਹ} ਮਲਾਹਾਂ ਦਾ ਬਹੁਤ ਆਉਣਾ-ਜਾਣਾ ਵੇਖਦੀ ਹੈ।
\VUsaut{6.0mm}
\VUfilet{22mm}
